% =========================================================
% Definition — Subspace
% =========================================================

\begin{tcolorbox}[colback=propbox, colframe=propborder, arc=2pt,
  left=6pt, right=6pt, top=4pt, bottom=4pt,
  title={\small\textbf{Definition (Subspace)}},
  fonttitle=\small\bfseries]
\begin{definition}[Subspace]
\label{def:subspace}
Let $V$ be a vector space over $\mathbf{F}$.
A subset $U \subseteq V$ is called a \textbf{subspace} of $V$ if $U$ is itself
a vector space over $\mathbf{F}$ with the same addition and scalar multiplication
as on $V$.
\end{definition}
\end{tcolorbox}

\begin{remark*}[Standard quantified statement]
\[
\begin{aligned}
&U \text{ is a subspace of } V \\
&\quad\Longleftrightarrow\quad
U \subseteq V
\text{ and } U \text{ is a vector space over } \mathbf{F} \\
&\qquad\text{under the operations inherited from } V.
\end{aligned}
\]
\end{remark*}

\begin{remark*}[Definition predicate reading]
\[
\operatorname{Subspace}(U,V).
\]
\end{remark*}

\begin{remark*}[Interpretation]
A subspace is a subset of a vector space that is closed under the vector-space
operations and therefore forms a vector space in its own right.
\end{remark*}

% =========================================================
% Theorem — Conditions for a Subspace
% =========================================================

\begin{tcolorbox}[colback=thmbox, colframe=thmborder, arc=2pt,
  left=6pt, right=6pt, top=4pt, bottom=4pt,
  title={\small\textbf{Theorem (Conditions for a Subspace)}},
  fonttitle=\small\bfseries]
\begin{theorem}[Conditions for a Subspace]
\label{thm:conditions-for-a-subspace}
Let $V$ be a vector space over $\mathbf{F}$ and let $U \subseteq V$.
Then $U$ is a subspace of $V$ if and only if the following three conditions hold:
\begin{enumerate}[label=(\roman*)]
  \item $0 \in U$;
  \item for all $u,v \in U$, one has $u+v \in U$;
  \item for all $a \in \mathbf{F}$ and all $u \in U$, one has $au \in U$.
\end{enumerate}
\end{theorem}
\end{tcolorbox}

\begin{remark*}[Standard quantified statement]
\[
\begin{aligned}
&U \text{ is a subspace of } V \\
&\quad\Longleftrightarrow\quad
0 \in U \\
&\qquad\land\;
\forall u,v \in U \;(u+v \in U) \\
&\qquad\land\;
\forall a \in \mathbf{F}\;\forall u \in U \;(au \in U).
\end{aligned}
\]
\end{remark*}

\begin{remark*}[Interpretation]
To check that a subset is a subspace, it is enough to verify that it contains
the zero vector and is closed under addition and scalar multiplication.
\end{remark*}

% =========================================================
% Definition — Sum of Subspaces
% =========================================================

\begin{tcolorbox}[colback=propbox, colframe=propborder, arc=2pt,
  left=6pt, right=6pt, top=4pt, bottom=4pt,
  title={\small\textbf{Definition (Sum of Subspaces)}},
  fonttitle=\small\bfseries]
\begin{definition}[Sum of Subspaces]
\label{def:sum-of-subspaces}
Let $V$ be a vector space over $\mathbf{F}$, and let
$U_1,\dots,U_m$ be subspaces of $V$.
The \textbf{sum} of $U_1,\dots,U_m$ is the set
\[
U_1+\cdots+U_m
=
\{u_1+\cdots+u_m : u_1 \in U_1,\dots,u_m \in U_m\}.
\]
\end{definition}
\end{tcolorbox}

\begin{remark*}[Standard quantified statement]
\[
\begin{aligned}
&v \in U_1+\cdots+U_m \\
&\quad\Longleftrightarrow\quad
\exists u_1 \in U_1 \;\cdots\; \exists u_m \in U_m
\text{ such that } v=u_1+\cdots+u_m.
\end{aligned}
\]
\end{remark*}

\begin{remark*}[Definition predicate reading]
\[
\operatorname{SumOfSubspaces}(W,\{U_1,\dots,U_m\},V)
\]
with
\[
W=U_1+\cdots+U_m.
\]
\end{remark*}

\begin{remark*}[Interpretation]
The sum of subspaces consists of all vectors that can be built by adding
together one vector from each subspace.
\end{remark*}

% =========================================================
% Theorem — Sum of Subspaces Is the Smallest Containing Subspace
% =========================================================

\begin{tcolorbox}[colback=thmbox, colframe=thmborder, arc=2pt,
  left=6pt, right=6pt, top=4pt, bottom=4pt,
  title={\small\textbf{Theorem (Sum of Subspaces Is the Smallest Containing Subspace)}},
  fonttitle=\small\bfseries]
\begin{theorem}[Sum of Subspaces Is the Smallest Containing Subspace]
\label{thm:sum-of-subspaces-smallest-containing}
Let $V$ be a vector space over $\mathbf{F}$, and let
$U_1,\dots,U_m$ be subspaces of $V$.
Then $U_1+\cdots+U_m$ is a subspace of $V$ containing each of
$U_1,\dots,U_m$.
Moreover, if $W$ is a subspace of $V$ containing each of
$U_1,\dots,U_m$, then
\[
U_1+\cdots+U_m \subseteq W.
\]
\end{theorem}
\end{tcolorbox}

\begin{remark*}[Standard quantified statement]
\[
\begin{aligned}
&U_1+\cdots+U_m \text{ is a subspace of } V, \\
&U_j \subseteq U_1+\cdots+U_m \quad \text{for } j=1,\dots,m, \\
&\forall W \subseteq V\;
\Bigl[
\bigl(W \text{ is a subspace of } V
\land U_1 \subseteq W \land \cdots \land U_m \subseteq W\bigr) \\
&\qquad\qquad\Rightarrow\;
U_1+\cdots+U_m \subseteq W
\Bigr].
\end{aligned}
\]
\end{remark*}

\begin{remark*}[Interpretation]
The sum is the smallest subspace that contains all the given subspaces.
It is the subspace generated by putting them together.
\end{remark*}

% =========================================================
% Definition — Direct Sum
% =========================================================

\begin{tcolorbox}[colback=propbox, colframe=propborder, arc=2pt,
  left=6pt, right=6pt, top=4pt, bottom=4pt,
  title={\small\textbf{Definition (Direct Sum)}},
  fonttitle=\small\bfseries]
\begin{definition}[Direct Sum]
\label{def:direct-sum}
Let $V$ be a vector space over $\mathbf{F}$, and let
$U_1,\dots,U_m$ be subspaces of $V$ such that
\[
V=U_1+\cdots+U_m.
\]
Then $V$ is called the \textbf{direct sum} of $U_1,\dots,U_m$ if each
vector $v \in V$ can be written in the form
\[
v=u_1+\cdots+u_m,
\qquad
u_1 \in U_1,\dots,u_m \in U_m,
\]
in exactly one way.
In this case, we write
\[
V=U_1\oplus\cdots\oplus U_m.
\]
\end{definition}
\end{tcolorbox}

\begin{remark*}[Standard quantified statement]
\[
\begin{aligned}
&V=U_1\oplus\cdots\oplus U_m \\
&\quad\Longleftrightarrow\quad
V=U_1+\cdots+U_m \\
&\qquad\land\;
\forall v \in V\;
\exists! \, u_1 \in U_1 \;\cdots\; \exists! \, u_m \in U_m
\text{ such that } v=u_1+\cdots+u_m.
\end{aligned}
\]
\end{remark*}

\begin{remark*}[Definition predicate reading]
\[
\operatorname{DirectSum}(V,\{U_1,\dots,U_m\},\mathbf{F}).
\]
\end{remark*}

\begin{remark*}[Interpretation]
A direct sum is a sum in which every vector has a unique decomposition
into pieces coming from the given subspaces.
\end{remark*}


\input{volume-iii/algebra/linear-algebra/notes/vector-spaces/figure-direct-sum}


% =========================================================
% Theorem — Condition for a Direct Sum
% =========================================================

\begin{tcolorbox}[colback=thmbox, colframe=thmborder, arc=2pt,
  left=6pt, right=6pt, top=4pt, bottom=4pt,
  title={\small\textbf{Theorem (Condition for a Direct Sum)}},
  fonttitle=\small\bfseries]
\begin{theorem}[Condition for a Direct Sum]
\label{thm:condition-for-direct-sum}
Let $V$ be a vector space over $\mathbf{F}$, and let
$U_1,\dots,U_m$ be subspaces of $V$ such that
\[
V=U_1+\cdots+U_m.
\]
Then the following are equivalent:
\begin{enumerate}[label=(\roman*)]
  \item \(V=U_1\oplus\cdots\oplus U_m\);
  \item if \(u_1 \in U_1,\dots,u_m \in U_m\) and
        \[
        u_1+\cdots+u_m=0,
        \]
        then
        \[
        u_1=\cdots=u_m=0.
        \]
\end{enumerate}
\end{theorem}
\end{tcolorbox}

\begin{remark*}[Standard quantified statement]
\[
\begin{aligned}
&V=U_1\oplus\cdots\oplus U_m \\
&\quad\Longleftrightarrow\quad
\forall u_1 \in U_1 \;\cdots\; \forall u_m \in U_m \\
&\qquad
\bigl(u_1+\cdots+u_m=0 \Rightarrow u_1=\cdots=u_m=0\bigr).
\end{aligned}
\]
\end{remark*}

\begin{remark*}[Interpretation]
A sum is direct exactly when the zero vector has only the trivial
decomposition as a sum of vectors from the given subspaces.
\end{remark*}

% =========================================================
% Theorem — Direct Sum of Two Subspaces
% =========================================================

\begin{tcolorbox}[colback=thmbox, colframe=thmborder, arc=2pt,
  left=6pt, right=6pt, top=4pt, bottom=4pt,
  title={\small\textbf{Theorem (Direct Sum of Two Subspaces)}},
  fonttitle=\small\bfseries]
\begin{theorem}[Direct Sum of Two Subspaces]
\label{thm:direct-sum-of-two-subspaces}
Let $U$ and $W$ be subspaces of a vector space $V$.
Then
\[
U+W=U\oplus W
\]
if and only if
\[
U \cap W = \{0\}.
\]
\end{theorem}
\end{tcolorbox}

\begin{remark*}[Standard quantified statement]
\[
\begin{aligned}
&U+W=U\oplus W \\
&\quad\Longleftrightarrow\quad
U \cap W = \{0\}.
\end{aligned}
\]
\end{remark*}

\begin{remark*}[Interpretation]
For two subspaces, directness is equivalent to saying that they overlap only
in the zero vector.
\end{remark*}








